\section*{الفصل \ref{ch:g10:measuring-the-earth} --- قياس الأرض}

\begin{solution}{exo:g10:measuring-the-earth:1}
تختفي السفينة بدءًا من هيكلها: فلو كان البحر مستويًا لصغرت فحسب وبقيت كلها
مرئية. وظل الأرض على القمر في الخسوف قوس من دائرة دائمًا: أما القرص المسطح
فيلقي ظلًّا مستطيلًا حين يُرى من حافته. وفي السفر شمالًا يعلو نجم القطب فوق
الأفق: ولو كانت الأرض مسطحة لبقي ارتفاعه ثابتًا تقريبًا.
\end{solution}

\begin{solution}{exo:g10:measuring-the-earth:2}
\omterm{def:g10:measuring-the-earth:geodesy}{الجيوديسيا} تقيس شكل الأرض وحجمها ومواقع النقاط عليها، وتغيرها مع
الزمن. \omterm{def:g10:measuring-the-earth:triangulation}{والتثليث} يحدد مواقع النقاط بسلسلة من المثلثات تُقاس كل زواياها؛
ولا يُقاس على الأرض إلا طول واحد، هو خط
القاعدة.
\end{solution}

\begin{solution}{exo:g10:measuring-the-earth:3}
$C = 50 \times 5000 = 250\,000$ ستاديون؛ $250\,000 \times \qty{157.2}{m} =
\qty{39300}{km}$.
\end{solution}

\begin{solution}{exo:g10:measuring-the-earth:4}
$111/60 \approx \qty{1.85}{km}$. وربع خط الزوال: $90 \times 111 \approx
\qty{10000}{km}$.
\end{solution}

\begin{solution}{exo:g10:measuring-the-earth:5}
$a - b = 6378.137 - 6356.752 = \qty{21.385}{km}$;
$f = 21.385/6378.137 \approx 0.003353 \approx 1/298$.
\end{solution}

\begin{solution}{exo:g10:measuring-the-earth:6}
في الجنوب: $\tan\theta_1 = 0.364$، $\theta_1 \approx \qty{20.0}{\degree}$. وفي الشمال:
$\tan\theta_2 = 0.456$، $\theta_2 \approx \qty{24.5}{\degree}$.
$\alpha = \qty{4.5}{\degree}$، إذن $C = (360/4.5) \times 500 = 80 \times 500 =
\qty{40000}{km}$.
\end{solution}

\begin{solution}{exo:g10:measuring-the-earth:7}
$C = 180 - 62 - 71 = \qty{47}{\degree}$.
$AC = 10.00 \sin 71^\circ/\sin 47^\circ = 10.00 \times 0.9455/0.7314 \approx
\qty{12.93}{km}$; $BC = 10.00 \sin 62^\circ/\sin 47^\circ \approx
\qty{12.07}{km}$.
\end{solution}

\begin{solution}{exo:g10:measuring-the-earth:8}
تغيّر الطول الجغرافي هو $10 + 20 = \qty{30}{\degree}$. وعند
\qty{60}{\degree} ش تبلغ الدرجة $111.3 \times 0.5 = \qty{55.65}{km}$، إذن
المسافة $30 \times 55.65 \approx \qty{1670}{km}$. وعلى خط الاستواء:
$30 \times 111.3 \approx \qty{3340}{km}$، أي ضعف المسافة.
\end{solution}

\begin{solution}{exo:g10:measuring-the-earth:9}
يلامس خط النظر البحر في $T$: والمماس لدائرة عمودي على نصف القطر في نقطة
التماس، فالزاوية عند $T$ قائمة. وحسب فيثاغورس: $(R+h)^2 = R^2 + d^2$، إذن $d^2 = 2Rh + h^2$؛ و$h$
ضئيل جدًا مقارنةً بالمقدار $R$، إذن $d \approx \sqrt{2Rh} = \sqrt{2 \times
6\,371\,000 \times 2} \approx \qty{5050}{m}$، أي نحو \qty{5}{km}.
\end{solution}

\begin{solution}{exo:g10:measuring-the-earth:10}
يكون المنحنى أشد انحدارًا في العروض الوسطى، من نحو \qty{30}{\degree} إلى
\qty{60}{\degree}، ومسطحًا قرب خط الاستواء والقطب. وعلى أرض مفلطحة يكون السطح
أقل انحناءً قرب القطبين؛ وكلما قل انحناء السطح، طالت المسافة التي يجب قطعها كي
يدور الرأسي (ومن ثم العرض الجغرافي) درجة
واحدة.
\end{solution}

\begin{solution}{exo:g10:measuring-the-earth:11}
$1072.2/9.646 \approx \qty{111.15}{km}$ للدرجة؛ $90 \times 111.15
\approx \qty{10004}{km}$، أي أكثر من ربع خط الزوال الحقيقي بمقدار \qty{2}{km}.
فمركز القوس قريب من \qty{46}{\degree}، حيث الدرجة أطول قليلًا من متوسط الدرجة
بين خط الاستواء والقطب
($10\,002.0/90 \approx \qty{111.13}{km}$).
\end{solution}

\begin{solution}{exo:g10:measuring-the-earth:12}
مع \qty{7.1}{\degree}: $C = (360/7.1) \times 5000 \approx 253\,500$ ستاديون
بدلًا من $250\,000$: أي أكبر مما ينبغي بنحو 1.4\,\% (فالزاوية الأصغر للمسافة
نفسها تعني أرضًا أكبر). أما المسافة، التي قدّرها سعاة يقطعونها مشيًا على طول
النيل، فكانت أقل دقة بكثير من الزاوية.
\end{solution}

\begin{solution}{exo:g10:measuring-the-earth:13}
في $ABC$: $C = \qty{50}{\degree}$، $AC = 6.000 \sin 75^\circ/\sin 50^\circ
\approx \qty{7.566}{km}$. وفي $ACD$: الزاوية عند $D$ هي $180 - 50 - 68 =
\qty{62}{\degree}$، و$AD = 7.566 \sin 68^\circ/\sin 62^\circ \approx
\qty{7.945}{km}$.
\end{solution}

\begin{solution}{exo:g10:measuring-the-earth:14}
$\cos\varphi = 1/2$ يعطي $\varphi = \qty{60}{\degree}$؛ و$\cos\varphi = 0.1$
يعطي $\varphi \approx \qty{84.3}{\degree}$. فدرجة الطول عند العرض $\varphi$
هي درجة خط الاستواء مضروبة في $\cos\varphi$: والنسبة لا تتوقف إلا على الزاوية،
لا على نصف القطر.
\end{solution}

\begin{solution}{exo:g10:measuring-the-earth:15}
عند خط الاستواء: $111.195 - 110.574 = \qty{0.621}{km} = \qty{621}{m}$؛ وعند
القطب: $111.694 - 111.195 = \qty{0.499}{km} = \qty{499}{m}$. فأخطاء من
مئات الأمتار في كل درجة، تتراكم درجة بعد درجة، أكبر بكثير من المتر: ولا بد أن
تعتمد المواقع المحوّلة من دوائر العرض وخطوط الطول على
الإهليلجي.
\end{solution}

\begin{solution}{pb:g10:measuring-the-earth:1}
\textbf{1.} $31.20 - 24.09 = \qty{7.11}{\degree}$.

\textbf{2.} $7.11/360 \approx 0.0198 = 1/50.6$، وهو قريب جدًا من الكسر
$1/50$ عند إراتوستينس.

\textbf{3.} $C = 788 \times 360/7.11 \approx \qty{39900}{km}$.

\textbf{4.} قوس خط الزوال بين دائرتي العرض: ففرق زاويتي الظهيرة هو فرق العرض
الجغرافي، الذي لا يقيس إلا البعد بين المكانين في اتجاه الشمال--الجنوب؛ ولا دور
لمركّبة المسافة في اتجاه الشرق--الغرب.

\textbf{5.} عند الظهيرة المحلية تكون الشمس في مستوي خط الزوال، وهو نفسه
للمكانين؛ وفي اليوم نفسه يكون لها الموقع نفسه بين النجوم، فلا تختلف الزاويتان
إلا بسبب انحناء الأرض.

\textbf{6.} $180 - 58 - 66 = \qty{56}{\degree}$.

\textbf{7.} $AC = 12.000 \sin 66^\circ/\sin 56^\circ = 12.000 \times
0.9135/0.8290 \approx \qty{13.223}{km}$.

\textbf{8.} $BC = 12.000 \sin 58^\circ/\sin 56^\circ \approx \qty{12.275}{km}$.

\textbf{9.} الزاوية عند $D$ هي $180 - 47 - 75 = \qty{58}{\degree}$؛
$AD = 13.223 \sin 75^\circ/\sin 58^\circ \approx \qty{15.061}{km}$.

\textbf{10.} $15.061 \cos 8^\circ \approx \qty{14.914}{km}$.

\textbf{11.} $14.914/111.13 \approx \qty{0.1342}{\degree}$، أي
$0.1342 \times 60 \approx 8.05$ دقيقة قوسية.

\textbf{12.} $111.51 - 110.57 = \qty{0.94}{km}$: \qty{940}{m}.

\textbf{13.} يدور الرأسي ببطء أكبر مع المسافة قرب القطب: فالسطح هناك أقل
انحناءً. إذن الأرض مفلطحة عند القطبين ومنتفخة عند خط
الاستواء.

\textbf{14.} $a - b = \qty{21.385}{km}$; $f = 21.385/6378.137 \approx 1/298$.

\textbf{15.} خيط الشاقول عمودي على سطح الإهليلجي. وفي القطع الناقص لا يمر
العمودي على المنحني بالمركز إلا عند خط الاستواء وعند القطبين؛ أما عند
\qty{45}{\degree} فيقطع المحور تحت المركز.

\textbf{16.} $10\,000\,000$ متر، أي \qty{10000}{km}.

\textbf{17.} $10\,001\,966/10\,000\,000 \approx \qty{1.0001966}{m}$.

\textbf{18.} أقصر مما ينبغي: فالجزء الحقيقي من عشرة ملايين يساوي $1.0001966$ من
متر 1799، إذن متر 1799 أقصر بنحو $1.97 \times 10^{-4}$ من
طوله، أي 0.02\,\%.

\textbf{19.} أقصر مما ينبغي. فالقضيب جزء من عشرة ملايين من ربع خط الزوال كما
حُسب في تسعينيات القرن الثامن عشر؛ وبما أن القضيب أقصر من جزء من عشرة ملايين
من ربع خط الزوال الحقيقي، فقد كان ربع خط الزوال المحسوب أقصر من
الحقيقي.

\textbf{20.} $1.0001966 - 1 = 0.0001966$ م: كان المتر الأول أقصر مما ينبغي
بنحو \textbf{\qty{0.2}{mm}} --- أي سماكة ورقتين تقريبًا على طول متر، وهي دقة
مدهشة لمسح أُجري بالمناظير والزوايا عبر بلد
بأكمله.
\end{solution}
